\begin{figure}[htbp]
\centering
\begin{tikzpicture}[font=\small]
  \tikzset{inputnode/.style={draw, circle, fill=black!5, minimum size=8mm},
           gate/.style={draw, rounded corners, fill=lightaccent, minimum width=15mm, minimum height=8mm, align=center},
           outputnode/.style={draw, circle, fill=warnbg, minimum size=9mm},
           wire/.style={->, >=stealth, thick}}

  \node[inputnode] (x) at (0,1.4) {$x$};
  \node[inputnode] (y) at (0,0.0) {$y$};
  \node[inputnode] (z) at (0,-1.4) {$z$};

  \node[gate] (notx) at (1.7,1.4) {NOT};
  \node[gate] (and1) at (3.5,0.9) {AND};
  \node[gate] (and2) at (3.5,-0.9) {AND};
  \node[gate] (or1) at (5.4,0.0) {OR};
  \node[outputnode] (f) at (7.0,0.0) {$f$};

  \draw[wire] (x) -- (notx);
  \draw[wire] (notx) -- (and1);
  \draw[wire] (y) -- (and1);
  \draw[wire] (x) to[out=-20,in=160] (and2);
  \draw[wire] (z) -- (and2);
  \draw[wire] (and1) -- (or1);
  \draw[wire] (and2) -- (or1);
  \draw[wire] (or1) -- (f);

  \node[align=center, font=\scriptsize] at (3.7,-2.0) {$f(x,y,z)=(\overline{x}\wedge y)\vee(x\wedge z)$};
\end{tikzpicture}
\caption{Egyszerű aciklikus logikai hálózat egy Boole-függvény kiszámítására}
\label{fig:zv10_logikai_halozat}
\end{figure}
